%%%%%%%%%%%%%%%%%%%%%%%%%%%%%%%%%%%%%%%%%%%%%%%%%%%%%%%%%%%%%%%%%%%%%%%%%%%%
\documentclass[a4j]{jarticle}

\usepackage{docmas_v2}
\usepackage{color}
\usepackage[dvipdfmx]{graphicx}
\usepackage{caption}
\captionsetup[figure]{
font = small,
}
%\usepackage{epsf}
%%%%%%%%%%%%%%%%%%%%%%%%%%%%%%%%%%%%%%%%%%%%%%%%%%%%%%%%%%%%%%%%%%%%%%%%%%%

\begin{document}

%insertfigure(ファイルの名前, 図タイトル, width)
\newcommand{\insertfigure}[3]{
\begin{figure}[htbp]
    \begin{center}
        \includegraphics[width={#3}]{#1}
        \caption[#2]{#2．}
        \label{fig:#2}
    \end{center}
\end{figure}
}

%insertfigure(ファイルの名前, 図タイトル, width，ref)
\newcommand{\insertreffigure}[4]{
\begin{figure}[htbp]
    \begin{center}
        \includegraphics[width={#3}]{#1}
        \caption[#2]{#2．#4より引用．}
        \label{fig:#2}
    \end{center}
\end{figure}
}

%inserttable(表タイトル, 表本体)
\newcommand{\inserttable}[2]{
\begin{table}[htbp]
    \begin{center}
        \caption{#1}
        \label{tb:#1}
        #2
    \end{center}
\end{table}
}

%inserteqnarray(式タイトル, 式本体)
\newcommand{\inserteqnarray}[2]{
\begin{eqnarray}
    \label{eqn:#1}
    #2
\end{eqnarray}
}


% 和文タイトル
\title{物語構造分析に基づくLLMを用いたインタラクティブ\\ストーリー生成システムによる創作支援}

% 英文タイトル
\etitle{Creation Support by Interactive Story Generation System Using LLM Based on Story Structure Analysis.}%DeepLにて翻訳 

% 著者名: 
%	・各著者を\quad（全角空白）区切りで列挙
% 	・著者名の直後に\afil{所属番号}を追加→所属番号を上付で出力（\textsuperscript{所属番号}と同じ）
% 	　複数機関へ所属している場合は番号をカンマ区切りで列挙（下記著者２参照）
% 　・Corresponding Authorについては所属の後に\thanksを続け，連絡先を記入
%	・英文著者はカンマ区切りで列挙

\author{渡邉謙吾\afil{1},%
	\thanks{連絡先：慶應義塾大学理工学部 \newline%
		　　　　　　神奈川県横浜市港北区日吉3-14-1\newline%
            　　　　　　kng.w@keio.jp}\quad%
	川村天\afil{2},
        小林伶央\afil{2},\
        有井知真\afil{1},\
        \\
        伊藤亮史\afil{1},\
        栗原聡\afil{1}\\
	Kengo Watanabe\afil{1}, Takashi Kawamura\afil{1}, Reo Kobayashi\afil{1}, Kazuma Arii\afil{1},\\ Akifumi Ito\afil{1}, and Satoshi Kurihara\afil{2} }


% 所属
\affiliation{%
	\afil{1} 慶應義塾大学理工学部\\ 
	\afil{1} Faculty of Science and Technology, Keio University\\
	\afil{2} 慶應義塾大学大学院理工学研究科\\ 
	\afil{2} Graduate School of Science and Technology, Keio University
}

\abstract{In recent years, the fluidity of media contents has expanded rapidly, leading to a growing demand for media contents. As a result, the challenges in creating stories have increased, and there is a growing need for story creation support. The development of AI technology, including large language models (LLMs), has been remarkable, and there are high expectations for these tools to aid in creative activities, though their utilization methods are not yet established. In particular, there is a demand for systems that offer meaningful support to individuals who are not familiar with AI. To address this, we have developed a web application that provides an intuitive interface for controlling LLMs based on narrative structure analysis, as well as interactive co-creation methods aimed at improving the quality of plot development. Our application includes features for generating plots based on the narrative structure analysis of Osamu Tezuka's manga "Black Jack" and functions to enhance plot quality through interactive interactions with the LLM, thus offering a collaborative co-creation approach. When utilized by professional story creators, our application received high evaluations for its usability and reflection of instructions, and it has been suggested as a useful interactive co-creation method. Through our study, we have explored approaches to co-creation between people and AI in support of story creation. }

\maketitle
\thispagestyle{empty}

\section{はじめに}
近年, ストリーミングサービスなどメディアコンテンツに関連するサービスが大きな広がりを見せており, 社会におけるメディアコンテンツの流動性が増している. 日々大量のコンテンツが消費され, 次から次へと新しいものが求められるようになった. 

このようなコンテンツ需要の高まりに応えるストーリー創作の困難は大きい. 多くのコンテンツは, ストーリーを骨格として形成され, その新奇性と多様性がコンテンツの魅力と供給に直結するためである. ストーリーの創作には高度な創造性を要し, コンテンツ需要が加速する中, 十分な新しさや多様性をもつ高品質なストーリーを, 社会の消費に見合う速度で生み出すことに大変な労力が必要になっているのである.

このような負担を軽減するため, 創造性を支援する技術が望まれており, AI技術の応用が期待されている.  

創造性とは, 価値のある新たな発想を生み出す能力のことであり, 価値には, 面白さや美しさ, 有益さといったさまざまな意味が含まれる\cite{ref:Boden09}. したがって, ストーリー創作における創造性の支援は, これまでにない面白さをもつ物語への気づきを与えるようなものであることが望まれる. 

従来のAIがこなせるタスクは単純なものに限られていたため, このような創造性を必要とするタスクに適用することは難しいという議論もあった\cite{ref:Boden09}. 

しかし現在, より高度で幅広いタスクにも対応可能な Foundation Model\cite{ref:Bommasani21}が登場し, AI の適用範囲は創造的活動にまで広がっていると注目されている\cite{ref:Frich19}. 特に, GPT-4\cite{ref:OpenAI23} をはじめとした大規模言語モデル(LLM) への期待は大きい. 

今日の LLM は, 自然言語処理におけるベンチマークにおいて驚異的な数値を叩き出す\cite{ref:Zhao23}. 自然言語を理解しているかのように応答し, 生成される文章は人が書いた文章に匹敵するほどである. 
%自然言語処理分野における様々なタスクで高い性能を発揮するようになっており\cite{ref:Zhao23}, ストーリー創作における創造性支援の可能性を大きく広げている.



様々な言語タスクに応用可能なLLMの有用さは, 創造性支援にも活用できると考えられ, これまでも, ストーリークリエーターの創造性支援に LLM を用いる試みが行われてきた\cite{ref:Mirowski23}.

しかしながら, その活用方法は未だ確立されていない. LLM から有用な出力を引き出すためにはプロンプトに工夫が必要であると知られており\cite{ref:Chen23}, AIの活用に馴染みのないストーリークリエーターでも, 負担なく有意義に LLM を活用することができる体系の構築が求められている. 

そこで我々は, 物語構造を活用し, プロのストーリークリエーターが LLM を用いたストーリー創作を効率的に行うことができる, ユーザーと LLM の仲介を担う御用聞き AI を web アプリとして開発した. このアプリケーションは画面の指示に従う直感的な操作で, 物語構造分析に基づいたプロット生成を実現する. さらに, 生成されたプロットの改善を目的とする, LLM との体系的でインタラクティブな共創手法を提供する. 我々はこのアプリケーションの開発により, LLM を活用したプロのストーリークリエーターの創作支援を探求した.

%そこで我々は, 物語構造を活用し, プロのストーリークリエーターが LLM を用いたストーリー創作を効率的に行うことができる web アプリケーションを開発した. ユーザーは, このアプリケーションの操作に複雑な入力を必要としない. 画面の指示に従い選択と入力を進めていくことで, 事前に分析された物語構造に沿うプロットの生成を行うことができる. さらに, このアプリケーションでは, 生成されたプロットの改善を目的とする, 体系的でインタラクティブな LLM とのやり取りを行うことができる. このアプリケーションの開発により, プロのストーリークリエーターの創作を支援する LLM の活用方法を模索した.

%研究の成果として、本研究チームは物語構造に基づくフレームワークを適用し、専門のクリエーターがLarge Language Models（LLM）を駆使してストーリーテリングを効率的に行うためのウェブベースのアプリケーションを開発した。このアプリケーションのユーザーインターフェースは複雑な操作を必要とせず、直感的に画面のガイダンスに従って選択肢を選び、入力を行うことで、事前に解析済みの物語構造を反映したプロット生成を容易に実現する。加えて、アプリケーションは生成されたプロットの洗練を促進する目的で、体系的かつ対話的なLLMとのインタラクションを可能にしている。このアプリケーションの導入により、プロフェッショナルなストーリークリエーターがLLMを活用する新たな方法論についての検討を行い、創造過程における支援を目指すものである。

\section{関連研究}
\subsection{創作活動におけるAIの活用}
AI技術の急速な発展に伴い, AIを創作活動に応用する研究が盛んに行われている. その分野は多岐に渡る. 

Ko ら\cite{ref:Ko23}は, ビジュアルアーティストがLTGM(Large-scale Text-to-image Generation Model)の導入についてどのような理解を持っているのかを調査し, LTGMを活用するユーザーインタフェースの設計に有用なデザインガイドラインを提示した. Louie ら\cite{ref:Louie20}は, AI の楽曲生成が非決定的に行われる問題に対処するため, 生成される楽曲を特定の声に制限し, 生成物の類似性を制御するセマンティックスライダーなどを実装したツールを提案した. Guzdial ら\cite{ref:Guzdial19}は, ゲームレベルデザイナーのパートナーとなるようなAIエージェント, Morai Makerを開発した.

このように, AI の創造的活動への適応は, 画像や音楽, ゲームデザインにまで広がっている.
\subsection{言語モデルとストーリー創作支援}
%言語モデルは, Trasformer\cite{ref:Vaswani17}の登場を起点とし, 著しい発展を遂げてきた. Transformer をベースとしたモデルには BERT\cite{ref:Devlin18} や GPT\cite{ref:Radford20} などがある. 学習データとモデルの大規模化により高い汎化性能を獲得し, GPTは, その系譜を GPT-4\cite{ref:OpenAI23}にまで進歩させている. 
言語モデルの急速な発展に伴い, これらをストーリークリエーターの創作支援に活用するシステムの研究が進められている.

大曽根ら\cite{ref:Osone21}は日本の小説家に用いられることを想定したAIシステムBunchoを開発した. Bunchoは, 日本語記事の大規模データセットと web 上の小説によって学習された GPT-2\cite{ref:Radford20} を利用しており, キーワードから, タイトルや比較的短いあらすじを生成することができる. Yuanら\cite{ref:Yuan22}は, LaMDA\cite{ref:Thoppilan22} と物語の共創が行えるインターフェース, Wordcraft を開発した. Wordcraft では, 自身で書いた文章の続きを書かせることや, 文章の一部を指定して, その部分の代替案を出力させることができる. これらのシステムは, 有意義な共創体験を与えるという評価を得た一方で, ストーリーに対する認識の欠如や生成物の一貫性が不足しているという評価を受けた.

%利用した作家からは, AIとの共創体験が楽しいものであったという評価を受けた. 一方で, 生成物の正確さや一貫性という観点では, 人の手で書かれたものよりも有意に低い評価を受けた.

 %AI との対話によって文章を改変していくことができるシステムであるが, AI の返答がうまく文脈を認識していないと感じられる課題が生じている. 

日笠ら\cite{ref:Higasa23}は, 物語構造を用いGPT-3\cite{ref:Brown20} を制御することで生成したプロットユニットを, 意味的類似度で繋ぎ合わせプロットの自動生成を行なった. ここでいうプロットとは, 物語の大きな流れを示す文章のことであり, プロットユニットとは, プロットを展開ごとに分割した要素のことを表す. この研究により, プロット生成における物語構造の有用性が示唆された. しかしながら, このシステムはクリエーターが直接操作できず, クリーエーターの意思を反映することができない.

%日笠ら\cite{ref:Higasa23}は, 物語構造を用いて GPT-3\cite{ref:Brown20} を制御し, プロットユニットを生成をすることで, プロットの自動生成に取り組んだ. ここでいうプロットとは, 物語の大きな流れを示す文章のことであり, プロットユニットとは, プロットを展開ごとに分割した要素のことを指す. 個別に生成したプロットユニットを意味的類似度で繋ぎ合わせ, 一つの大きなプロットを作り上げた. このことにより, プロット生成において, 物語構造の活用が有用であることを示唆した. しかしながら, このシステムはストーリークリエーターが自身の手で動かすことのできるシステムではなく, クリーエーターの意思を反映することができない.

\section{提案システム}
\begin{figure}[htbp]
\begin{center}
\includegraphics[width=40mm]{img/3/system.png}
\caption{提案システムのイメージ図.}
\label{system}
\end{center}
\end{figure}
\vspace{-1.5zh}
%提案手法では, アプリケーションの機能を一つ一つきちんと説明する
我々は, プロのストーリークリエーターを対象とした, ストーリー創作支援を行う web アプリケーションを開発した. 本アプリケーションは, 手塚治虫の漫画『ブラック・ジャック』\cite{ref:Tezuka77}における物語構造分析や, キャラクター分析により得られたデータを LLM の制御に活用する. 機能は大きく2つのセクションに分かれ, 一つはプロットの生成を行うセクション, もう一つはプロットの編集を行うセクションである. 分析に基づく『ブラック・ジャック』の新作プロットの生成を, 直感的な操作により実現するインターフェースと, プロットの改善を目的とした LLM とのインタラクション手法を提供する. いずれのセクションにおいても, ストーリークリエーターと LLM との効果的なやりとりを促すフレームワークを実装した.

%インターフェースを通じた直感的な操作により, 分析に基づく『ブラック・ジャック』の新作プロットの生成を可能とし, プロットの質改善を目的とした LLM とのインタラクション手法を提供する.

%分析に基づいた『ブラック・ジャック』の新作プロットの生成を, インターフェースを通じた直感的な操作により実現し, LLM との効果的なインタラクション手法を提供する.

%我々は, プロのストーリークリエーターを対象とした, ストーリー創作支援を行う web アプリケーションを開発した. 本アプリケーションは, 手塚治虫の漫画『ブラック・ジャック』\cite{ref:Tezuka77}における物語構造分析や, キャラクター分析により得られデータを LLM の制御に活用する. ユーザーは本アプリケーションのインターフェースを通じて, 直感的な操作により, 分析に基づく『ブラック・ジャック』の新作プロットを生成するこができる.

%また, 本アプリケーションは, プロットの改善を目的とする LLM とのインタラクティブな共創手法を提供し, クリエーターのストーリー創作を支援する.

%我々は, プロのストーリークリエーターを対象とした, ストーリー創作支援を行う web アプリケーションを開発した. 本アプリケーションは, 手塚治虫の漫画『ブラック・ジャック』\cite{ref:Tezuka77}における物語構造分析や, キャラクター分析により得られデータを LLM の制御に活用し, クリエーターの物語生成過程を容易にする. ユーザーは, 本アプリケーションのインターフェースを通じて, 直感的な操作により, 分析に基づく『ブラック・ジャック』の新作のプロットを生成するこができる.  また, 本アプリケーションは, プロットの改善を目的とする LLM とのインタラクティブな共創手法を提供し, クリエーターのストーリー創作を支援する. 

%直感的な操作を通じて, 分析に基づく『ブラック・ジャック』の新作のプロットを容易に生成することができる. また, 本アプリケーションは, プロットの改善を目的とする LLM とのインタラクティブな共創手法を提供し, クリエーターのストーリー創作を支援する. 

%我々は, プロのストーリークリエーターのストーリー創作を支援する web アプリケーションを開発した. このアプリケーションは, 手塚治虫原作の漫画『ブラック・ジャック』[手塚治虫 77]に対して行われた物語構造分析や, キャラクターの分析データを活用し, クリエーターの LLM を用いたストーリー生成を容易にする. ユーザーは, このアプリケーションを用いることで, 簡単な選択と入力により, 分析に基づいた『ブラック・ジャック』の新作プロットを生成することができる. また, このアプリケーションは, プロットの改善を目的とする LLM とのインタラクティブな共創手法を提供し, クリエーターのストーリー創作を支援する. 

%機能は大きく2つのセクションに分かれ, 1つはプロットの生成を行うセクション, もう一つはプロットの編集を行うセクションである. いずれのセクションにおいても, ストーリークリエーターと LLM との効果的なやりとりを促すフレームワークを提供する. 

本アプリケーションに接続した LLM は GPT-4\cite{ref:OpenAI23} である. GPT-4 は, OpenAI 社が開発した LLM であり, あらゆる自然言語処理タスクにおいて高い性能を発揮する\cite{ref:OpenAI23}. API も公開されていることから, インターフェースへの接続も容易であり, 本実験では GPT-4 を用いることとした.

\subsection{物語構造}
物語構造とは, 物語が進行するパターンを構造化した概念を指す. この研究では, 約300種のロシアの魔法昔話を31の要素の組み合わせにより記述できると主張したProppの研究\cite{ref:Propp64} などがよく知られている. 

%物語構造とは, 物語が進行するパターンを構造化した概念を指し, 物語の構造化は, 物語の特性や構成を対象とする学問分野, 物語論(narratology)の中で取り扱われてきた. この分野における研究では, 約300種のロシアの魔法昔話を31の要素の組み合わせにより記述できると主張したProppの研究\cite{ref:Propp64} などがよく知られている. 

物語の自動生成に関する研究において, 物語構造を利用した研究はこれまでもさまざま行われてきた. Proppが見出した31機能を並べ替え, プロットの自動生成を試みる研究\cite{ref:Jaya11}などである. しかしながら, Proppが見出した31の要素はあくまで特定のジャンル, ロシアの魔法昔話の分析によって得られたものであり, 一般的な物語の作成に利用することはできないと考えられる. 

そこで村井\cite{ref:Murai22}は, 特定ジャンルの物語に限定されない, 一般的な物語の展開パターンを抽出する目的として, 『ブラック・ジャック』を分析対象とし, 物語構造分析を行った. 村井の研究によると『ブラック・ジャック』は, 複数ジャンルの物語パターンを含む短編集であり, 分析可能な話数も十分である. 加えて作品自体の評価も高い. 一般的な物語の展開パターンを抽出するにあたり, 異なる作者, 異なる作品の様々なジャンルの物語を組み合わせて分析することも考えられるが, しかしその場合, 作者や作品の違いによって, 演出上の表現技法などが大きく異なることとなり, 分析が困難になると考えられる. このような理由により, 同一作者, 同一作品内において, 複数ジャンルの物語を含み, 国内外で高い評価を持つ『ブラック・ジャック』は一般的な物語構造の分析に適した作品であるといえる\cite{ref:Murai22}.

そこで本研究では, 村井の物語構造分析により得られた『ブラック・ジャック』の物語構造, 展開パターンを LLM の制御に用いることとし, 『ブラック・ジャック』の新作プロットの創作支援に取り組んだ. 

『ブラック・ジャック』の新作プロットの創作支援により得た知見は, 同作品が複数ジャンルの物語構造を含む複合的な作品であるため, 一般的な物語作品を生成するための汎用的な手法に貢献すると考えられる.

\subsection{プロットの生成}
プロット生成セクションでは, 図\ref{flow}に示すフローに沿って操作を進める.
\begin{figure}[htbp]
\begin{center}
\includegraphics[width=65mm]{img/3/flow.png}
\caption{プロット生成セクションにおけるフロー.}
\label{flow}
\end{center}
\end{figure}

フローに沿って進めることで, 『ブラック・ジャック』の分析に基づくプロット生成を行うことができる. どのステップにおいても操作は簡潔で, 画面の指示に従い選択と入力を行うことにより容易に次のステップへと移行できるようになっている.

ここでいうプロットとは, 物語の全体のおおまかな流れを表す文章のことであり, 本アプリケーションで生成することができるプロットは, 1000字から 3000字程度のプロットである.

物語そのものではなくプロットを生成する理由は二つある. 一つは, 物語そのものには, 構造の上に演出上のカスタマイズがなされるため, 物語構造による制御の効果を直接活かせないと考えられること. もう一つは, 本アプリケーションがクリエーターの代替となるものではなく, クリエーターの創作支援を目的としたものであるため, 創造性を掻き立てる形式としてプロットの形式を選択したことである.
\subsubsection{ジャンル・テーマの設定}
ここでは, 物語のジャンルおよびテーマを, 『ブラック・ジャック』の分析により得た, 「SF」や「人間ドラマ」など, ジャンル64種と, 「親子愛」や「友情」など, テーマに関するキーワード152種の中から複数選び設定することができる.

分析に基づいた選択肢をあらかじめ用意しておくことで, 創作の手間を削減し, 多岐にわたるキーワードの提示により, 創造性を刺激する. 

選択は, 自身の意思により行うこともできれば, ランダム選択機能を用いて行うこともできる. 選び取る時間の削減とランダムネスによる創造性支援を行う\cite{ref:Osone21}.

\subsubsection{象徴表現の生成}
ここでは, 選択したテーマや自身で考えたテーマの象徴表現を生成することができる. 

ここで言う象徴表現とは, 鳩で平和を表すように, 抽象的な概念を別のものに置き換え表す表現技法のことである. この表現技法は, 小説や映画, 詩など様々な場面で用いられ, 物語に何層もの意味を加えて作品を豊かにするのに役立つものである\cite{ref:Norva23}.

象徴表現を生成することで, プロットに深みを与える手段や気づき, ヒントを獲得することができる.

象徴表現の生成例:

\begin{table}[hbtp]
\small
\begin{center}
{
\label{adopt_t}
\begin{tabular}{cc} 
信頼  →& 手を繋ぐ行為
\end{tabular}
}
\end{center}
\end{table}

\vspace{-3zh}

\subsubsection{テーマの言い換えの生成}
ユーザーは, テーマの言い換えの生成を行うことで, 単語形式で設定したテーマを文章の形式で獲得することができる. 
% \begin{table}[hbtp]
% \begin{center}
% {\small
% \caption{言い換えの生成例}
% \label{themachange}
% \begin{tabular}{cc}\hline
% テーマ & 言い換え\\ \hline
% 親子愛  &\begin{tabular}{c}
% 親と子が互いの理解や共感を深め、\\困難を乗り越える過程を描く物語
% \end{tabular} \\ \hline
% \end{tabular}
% }
% \end{center}
% \end{table}
% \begin{center}
% \begin{minipage}[htbp]{70mm}
% \small
% 「親子愛」 → 「親と子が互いの理解や共感を深め、困難を乗り越える過程を描く物語」
% \end{minipage}
% \end{center}

LLM は物語を生成する際, 入力として与えた単語をそのままの形で登場させてしまう傾向がある. しかし, 物語のテーマは, 作品全体を通して伝わることが望ましく, 生成するプロットの中に直接単語として現れてしまうことは好ましくない.

テーマの言い換えの生成により, 文章形式でテーマを設定し直すことで, 選択したキーワードが生成するプロットの中に直接出てくることを防ぐことができる.

テーマの言い換えとして生成された文章は, ユーザー自身の手で書き換えることもでき, LLM とクリエーターとの共創を実現する.

言い換えの生成例:
\begin{table}[hbtp]
\begin{center}
{\footnotesize
\label{adopt_t}
\begin{tabular}{cc} 
親子愛  →&\begin{tabular}{c}
親と子が互いの理解や共感を深め、\\困難を乗り越える過程を描く物語\end{tabular} 
\end{tabular}
}
\end{center}
\end{table}
\vspace{-1.5zh}

なお, 同じキーワードに対する言い換えの文章が同じ文章になるとは限らない. 

\subsubsection{キャラクターの設定}
ここでは, 物語に登場させる, 『ブラック・ジャック』の既存キャラクターと, 自身で作り上げる新規キャラクターの設定を行うことができる.

既存キャラクターには, 「ブラック・ジャック」や「ピノコ」など, 漫画『ブラック・ジャック』に登場する主要なキャラクターを選択肢として用意しており, ユーザーはその中からプロットに登場させるキャラクターを選ぶことができる. それぞれのキャラクターに対し, 「主人公」や「協力者」など物語上での役割や性格といった属性を, 分析に基づいた選択肢の中から設定することができ, 既存キャラクターの言動が原作からかけ離れたものにならぬよう工夫を施している.

\begin{figure}[htbp]
\begin{center}
\includegraphics[width=81mm]{img/3/chara.png}
\caption{新規キャラクターの設定画面.}
\label{chara}
\end{center}
\end{figure}

新規キャラクターには, 図\ref{chara}のように, 名前, 役割, 性別, 年齢, 身分・職業, 特徴, 性格という属性を設定することができる. それぞれの項目について『ブラック・ジャック』のキャラクター分析に基づいた選択肢が用意され, その選択肢の中から属性を選ぶことも, 自身で思いついた言葉を当てはめることもできる. ランダム選択機能も備えており, 設定の負担を軽減する.

登場させたい人数分, 新規キャラクターの属性を設定し終えると, 図\ref{chara}に示すように, それぞれのキャラクターに関するプロフィールを生成することができる. ここで言うプロフィールとは, キャラクターの背景や人物像をより詳細に説明する文章のことであり, これは作品上におけるキャラクターの言動に影響する\cite{ref:Mirowski23}. プロフィールは, 設定された属性をもとに生成され, 自身の手で自由に書き換えることもできる. 

以上のような機能により, キャラクター創作を支援する.

\subsubsection{物語展開因子と極性の設定}
物語展開因子とは, 物語の展開に影響を与えるとされる要素のこと\cite{ref:Murai21}である. また, ここで言う極性とは, 物語それぞれの場面における雰囲気を表す.

ここでは, 『ブラック・ジャック』の分析により得られた, 「犯罪と罪の報い」や「夫婦」などの物語展開因子57種の中から, 含めたい要素を複数選択することができる. また, 物語の序盤, 終盤それぞれの極性を, -1(negative), 0(natural), 1(positive)の3つの数字の中から選び設定することができる. 例えば, 序盤に -1, 終盤に 1 を選ぶと, 物語は暗い雰囲気から始まり, 最後には明るい雰囲気で終わるようにすることができる.

\subsubsection{物語構造の設定}
ここでは, 『ブラック・ジャック』の物語構造が, 前のステップで選択された物語展開因子および極性に基づき, 4 から 11 の展開数ごとに複数提案される. ユーザーはその中から生成するプロットの物語構造を一つ設定することができる. これにより, 自身が含めたい要素と適合した, 分析に基づく物語構造をプロットの骨格として定めることができる.
%ここでは, 前のステップで選択された物語展開因子および極性に基づき提案される『ブラック・ジャック』の物語構造の中から, プロットの構造を一つ定めることができる. 
%ここでは, 生成するプロットの物語構造を, 前のステップで選択した物語展開因子および極性に基づき提案される『ブラック・ジャック』の物語構造から一つ選択することができる.

例えば, 提案される物語構造には, 展開数5の場合, 「探索」→「依頼」→「看病・治療」→「情報開示」→「真相発覚」などがある.

\subsubsection{自由指示・時代・場所の設定}
ここでは, 物語の舞台となる時代や場所について設定することができ, また, これまで定めてきた設定を確認することができる. これまでの設定を変更したり, 新たに追加したい要素を自由記述で設定することもできる. これまでの過程で受けた創造性の刺激を自身の言葉でアウトプットする場にもなっている.

\subsubsection{タイトルとあらすじの生成}
ここでは, これまで設定してきた情報をもとに, 1文から2文程度のあらすじとそれに付随するタイトルをそれぞれ10本生成することができる. 

ここで言うあらすじとは, 生成するプロットの元となる短い文章のことを指す.

タイトルのような物語全体に反映するキーワードはクリエーターの創造性を刺激する. また, 1文から2文程度のあらすじは, 物語の方向性に関するイメージを与え, 創造性の支援になると考えられる.

ユーザーは, 生成されたタイトルおよびあらすじから, プロットに活用するものをそのままの文章で選択することも, その内容を自身で加筆修正することもできる. この機能により共創支援を提供する.

\subsubsection{アイデアの生成}
ここでは, これまでの設定をもとに, 具体的な出来事や道具など, プロットの中で用いることができそうな物語全般に関するアイデアを生成することができる. 生成されたアイデアは, これまでの情報をプロットにまで拡張する材料になり得る.

\subsubsection{シーンごとの主観視点等最終設定およびプロットの生成}
ここでは, プロットの各シーンを誰の視点で描くか設定するとともに, これまで定めてきた設定を確認, 修正することができる.
\begin{figure}[htbp]
\begin{center}
\includegraphics[width=40mm]{img/3/settings.png}
\caption{設定の最終確認画面.}
\label{settings}
\end{center}
\end{figure}

シーンごとに主観視点を決められるこの機能は, 本アプリケーションの制作段階で, クリエーターから頂いた意見をもとに実装した機能であり, シーンごとに視点を変えることで, 物語に複雑性を持たせる狙いがある.

ここでの確認, 修正がプロット生成を行う前の最終段階であり, ここで確定させた設定をもとにプロット生成が行われる.

\subsection{プロットの編集}
本アプリケーションでは, プロット生成セクションの実装にとどまらず, LLM とのインタラクティブなやり取りにより, プロットを反復的に洗練させていくことができるプロット編集セクションを実装した. このセクションでは, 既存のプロットを, 新たな指示や要素を加え新しいプロットに書き換えることができ, 新旧プロットのより満足のいく方を採用していくことができる. この過程を反復して行うことでプロットの質を高めることができる. この機能によりクリエーターと LLM とのプロット生成における共創支援を実現する.

%本アプリケーションは, プロットを生成するだけでなく, プロットの改良を目的とした LLM とのインタラクティブなやり取りを可能にするプロット編集セクションを設けている. プロット編集を実行することで, 既存のプロットに, 新たな指示, 要素を加えた, 新たなプロットを生成させることができる. 新しく生成されたプロットと, 既存のプロットを比較し, どちらが良いか選択していくことで, プロットをより良いものにしていくことができる.
実装したプロット編集機能の一覧を以下に示す.
\begin{itemize}
\setlength{\itemsep}{0cm}
\item フリーワードによるやり取り
\item プロット中の単語・文章に関する質問
\item アイデア生成
\item 伏線の処理
\item 一部シーン変更
\item 全体的な雰囲気の変更
\item シーンの入れ替え
\item シーンごとの視点の変更
\item 登場人物の追加
\item 登場人物の役割変更
\item プロットの部分修正
\item 不自然な日本語の修正
\item プロット1シーンの起承転結への分割
\item 手作業での編集
\end{itemize}

\begin{figure}[htbp]
\begin{center}
\includegraphics[width=55mm]{img/3/flow2.png}
\caption{プロット編集セクションにおけるフロー.}
\label{flow2}
\end{center}
\end{figure}

以下では, これらの機能のうち, 実験でよく使われた3つの機能について説明する.

\subsubsection{フリーワードによるやり取り}
この機能を用いることで, ユーザーは生成されたプロットについて自由記述により LLM とやり取りを行うことができる. 自身の言葉でプロットに対する質問を投げかけることや, 新たな指示によりプロットを書き換えることができる.
% \begin{figure}[htbp]
% \begin{center}
% \includegraphics[width=70mm]{img/3/freequestion.png}
% \caption{フリーワードによるやり取り画面.}
% \label{free}
% \end{center}
% \end{figure} 

\subsubsection{プロット中の単語・文章に関する質問}
この機能では, プロットの中における, 説明が不十分であると思われる単語や内容をより深掘りしたいと思う文章について, 詳細な説明と解釈の例を回答として得ることができる. 
\begin{figure}[htbp]
\begin{center}
\includegraphics[width=81mm]{img/3/question.png}
\caption{プロット中の単語・文章に関する質問機能実行の様子. (1): 内容をより詳しく知りたいプロット上の単語や文章をテキストホルダーに入力する. (2): 詳しく聞くを押すと(1)で入力したものを深掘りした文章が生成される. (3): 回答にチェックをつける. (4): プロットに組み込むを押すと(3)でチェックをつけた回答がプロットに反映される.}
\label{question}
\end{center}
\end{figure} 

単に質問を投げかけるだけでは, プロットに書かれていること以上の内容が回答に含まれないため, プロットの範囲外のことも含め回答がなされるよう工夫を施した.

得られた回答を用いて新たなプロットを生成させることもでき, プロットを改善していくことができる.

\subsubsection{一部シーン変更}
この機能は, 選んだシーンの内容を変更することができる機能である. 他のシーンを変えることなく選択したシーンの内容のみを変更することができ, 変更されたシーンが, 矛盾なく前後のシーンと繋がり, プロットとしての一貫性が保たれるよう工夫を施した. 操作は単純で, 変更したいシーンの横にあるチェックボックスにチェックを入れるだけで良い. 操作の簡便性によりクリエーターの負担を軽減する. 

\subsubsection{新旧プロットの比較}
プロット編集セクションでは, 編集機能を用いてプロットの書き換えを行うたび, 新旧のプロットを比較することができる. 新旧プロットは画面上左右に表示され, それらを比較した結果として, 既存のプロットを採用するか, 新しいプロットを採用するか, 2つのプロットを手作業で統合するか選ぶことができる. 繰り返し編集機能を用いることでプロットを満足のいくものへと改良していくことができる.


\section{具体的な活用}
本アプリケーションの実践的な有用性を調査する目的として, プロのストーリークリエーター数グループに, 商業誌に掲載される『ブラック・ジャック』の新作\footnote{TEZUKA2023}のプロットを, 本アプリケーションを用いて制作してもらった.
%本研究では, 開発したアプリケーションの実践的な有用性を調査する目的として, プロのストーリークリエーターに, 出版誌に掲載する『ブラック・ジャック』の新作[TEZUKA 2023]のプロットを, 本アプリケーションを用いて制作してもらう実験を実施した.

プロットの制作期間は, 2023年07月08日から2023年08月02日までのおよそ1ヶ月間.

参加した被験者グループは5グループで, グループAは30代女性1名, グループBは60代男性1名, グループCは40代女性1名と60代男性1名, グループDは60代男性1名, グループEは女性1名と60代男性1名という構成であった.

それぞれのグループにおける創作活動歴と ChatGPT\cite{ref:OpenAI23}等, LLM の使用頻度は表\ref{group}に示す通り. 各被験者グループ内で最も大きな値を示している.

\begin{table}[hbtp]
\begin{center}
{\caption{被験者グループの構成}
\label{group}
\footnotesize
\begin{tabular}{ccc} \hline
グループ & 創作活動歴 & LLMの使用頻度  \\ \hline
A & 14年 & 月に数回使用する\\
B & 37年 & 全く使用したことがない\\
C & 24年 & ほとんど使用しない\\
D & 42年 & 全く使用したことがない\\
E & - & 全く使用したことがない \\
\hline
\end{tabular}
}
\end{center}
\end{table}
\

本実験では, 各被験者グループのプロット制作活動におけるアプリケーションの操作ログを取得した. さらに, 生成されたプロットの一貫性やアプリケーションの操作性, アプリケーションを活用したことによる創造性への刺激などについて, 5段階のリッカート尺度および自由記述により評価を求めるアンケートを実施した.

\section{結果および考察}

\subsection{アンケート結果および考察}

\subsubsection{生成されたプロットに対する満足度}
\begin{figure}[htbp]
\begin{center}
\includegraphics[width=70mm]{img/5/quality2.png}
\caption{本アプリケーションが生成するプロットの各項目における満足度を表すアンケート結果.}
\label{quality}
\end{center}
\end{figure}
図\ref{quality}より, 本アプリケーションは, 目新しさの提供に課題がある一方で, 幅広さや面白さといった観点で, プロのストーリークリエーターから見ても創作活動に役立つものであったことがわかる.

複合的な作品である『ブラック・ジャック』の分析データを, プロットの生成に活用したことが, LLM に幅広い言葉の出力を促したのではないかと考えられる.

一方, 目新しさの課題については, 出力が学習データに制限される LLM を創造的活動に用いる上で重要な課題の一つであるといえる.

\subsubsection{LLM との共創における貢献度}
\begin{figure}[htbp]
\begin{center}
\includegraphics[width=70mm]{img/5/creativity_control.png}
\caption{本アプリケーションにおける LLM との共創を促す貢献について評価されたアンケート結果.}
\label{qc}
\end{center}
\end{figure}
図\ref{qc}より, 作業時間の短縮に関する貢献を除き, 低い評価を受けていないことがわかる. 特に, 使いやすさと指示の反映に関する項目について高い評価を受けており, 本アプリケーションが LLM との創作活動を向上させる可能性を示唆した結果であるといえる.

作業時間については, LLM からのレスポンスの遅さに不満の声が聞かれ, より高速に応答する LLM がこの問題を解決すると考えられる.

創造性への刺激については, まずまずの評価を得ているが, 目新しさを提供する課題の取り組みにより, さらなる貢献をもたらすことができると考える.

\subsection{ログの結果および考察}

\subsubsection{作業時間}
プロット制作という創造的活動の作業時間を正確に把握することは難しい. 実際に手を動かしていなくとも, アイデアや構想を練ることなどに時間は使われるからである. 一方で創作活動の時間効率性を高めることは, 創作支援の重要な項目の一つであると考えられる. よって本論文では, 本アプリケーションの操作ログより作業時間を概算し, 本アプリケーションの時間効率への寄与について考察する.
\begin{figure}[htbp]
\includegraphics[width=60mm]{img/5/time_bar_A1A2.png}
\caption{ボタンの押下間隔から算出した作業時間. 被験者グループAは制作期間内において2つのプロットを制作したため作品ごとに作業時間を集計した. Average1 は全体平均, Average2 は外れ値を除いた平均値を表す.}
\label{time}
\end{figure}

作業時間は, ログデータとして集計した, 本アプリケーション画面のボタン押下間隔のうち, 98 パーセンタイルから 99 パーセンタイルの間に位置する 1 時間を区切りとして算出した. つまり, ボタン押下間隔が 1 時間以上開いている値は作業時間に含めずに, ボタン押下間隔が 1 時間未満の値を全て足し上げ算出した値を便宜的に作業時間とみなすこととした.

図\ref{qc}で示した, 「作業時間が短くなったか」 というアンケートへの回答のうち「短くなったとは思わない」と回答したのは, グループ C と E であり, そのほかのグループは「短くなったと思う」と回答している. 図\ref{time}からわかる通り, グループ C は作業時間を最も費やしたグループである一方, グループ E は最も作業時間の少ないグループであった.

グループ E が作業時間の短縮を感じなかったのは, LLM の応答の遅さや, 自身のプロット制作にかかる時間が元々十分に早いことなどが理由であると考えられる. 

\begin{table}[hbtp]
\small
\begin{center}
{\caption{各グループの編集機能の利用数}
\label{edit_action}
\begin{tabular}{ccccc} \hline
A & B & C & D & E  \\ \hline
116 & 60 & 176 & 33 & 45\\
\hline
\end{tabular}
}
\end{center}
\end{table}
\vspace{-1.5zh}

表\ref{edit_action}に示した通り, グループ C は最も多く編集機能を用いたグループでもある. これより, 編集機能における操作性などが時間効率に影響していると考えられる.

一方で, 作業時間の短縮を感じたと回答したグループはいずれも作業時間が 10 時間未満であり, このことに本アプリケーションの時間効率への寄与があったと考えられる.

\subsubsection{編集機能により生成したプロットの採用率}


\begin{table}[hbtp]
\small
\begin{center}
{\caption{グループごとの編集機能により生成したプロットの採用率とそれらの平均}
\label{adopt_t}
\begin{tabular}{cccccc} \hline
A & B & C & D & E & mean \\ \hline
0.71 & 0.91 & 0.70 & 0.50 & 0.63 & 0.69\\
\hline
\end{tabular}
}
\end{center}
\end{table}
表\ref{adopt_t}より, 編集機能を用いて, 既存のプロットを新たなプロットへと書き換えた際, ほとんどのグループで新旧プロットのうち新しいプロットを採用する割合が大きいことがわかる. 各グループの採用率の平均もおよそ 70 \% と高い割合を示している. 

この結果は, 表\ref{qc}で示したアンケート結果, 「編集機能により面白さは増加したか」という問において高い評価が得られた結果と符合している. つまり, これらの結果は, LLM とのインタラクティブなプロット創作を可能とする本アプリケーションの編集機能が, プロットの質向上に寄与することを示唆する結果であるといえる.

% \begin{figure}[htbp]
% \begin{center}
% \includegraphics[width=70mm]{img/5/comparison.png}
% \caption{編集機能により生成したプロットの採用・不採用数}
% \label{adopt}
% \end{center}
% \end{figure}

\subsubsection{用いられた編集機能の割合}

\begin{figure}[htbp]
\includegraphics[width=65mm]{img/5/edit_action_pie_chart_N.png}
\caption{グループごとに正規化した各編集機能の利用率}
\label{edit_action}
\end{figure}

図\ref{edit_action}より, 編集機能の中でフリーワードでやり取りする機能が最も多く利用され, グループごとに正規化した利用率の中でも50\%以上を占めることがわかる. 

この結果から, 生成されたプロットが, 型にはまらない自由な発想を引き出した可能性や, プロのストーリークリエーターが自身の言葉による制御を望んていたこと, プロットの多様性により, 機能を限定しない編集が必要となっていたことなどが考えられる.

フリーワードでのやり取りに対する支援の拡充が求められていると考える.

\section{結び}
本論文において我々は, ストーリークリエーターの創作を支援する web アプリケーションを提案した. 本アプリケーションは, 物語構造分析に基づく LLM の活用を, 直感的な操作により実現するとともに, LLM とのインタラクティブなプロット創作手法を提供する. 商業誌に掲載される作品のプロット制作に利用してもらい, 本アプリケーションが高い操作性と満足のいく指示の反映をもたらし, プロットの質向上を目的とする LLM との共創手法に有用である可能性が示唆された. 今後は使用者の範囲を拡大した実証実験を実施し, より信頼性のある有用性の確認を行う. 

本研究がストーリー創作支援におけるクリエーターと LLM との共創のあり方に貢献することを期待する.
%私たちは、BunChoという日本語版GPT-2を使い、キーワードをタイピングすることで小説のタイトルやあらすじを生成できる日本の小説共創AIインターフェースを提案します。BunChoは利用者がTRPGに似た対話形式で創造的な小説を共創するのを支援します。私たちのユーザースタディでは、作家自身による書かれたあらすじとBunChoにより共創されたあらすじを比較しました。その結果、利用者は自分自身だけで作成するよりもBunChoと共創することをより楽しんでいることが示されました。私たちは、AIとの共創による日本の物語作りの未来の一つの方向性を示したと信じています。さらに、私たちはBunChoの共創におけるAI TRPGの効果とユーザーエクスペリエンスを評価する予定です。

\section*{謝辞}
本研究はNEDO・人と共に進化する次世代人工知能に関する技術開発事業「インタラクティブなストーリー型コンテンツ創作支援基盤の開発」の助成を受けたものである．

%%%%%%%%%%%%%%%%%%%%%%%%%%%%%%%%%%%%%%

\begin{thebibliography}{99}
\small

\bibitem{ref:Boden09}
Boden, M. A.: Computer Models of Creativity, \textit{AI Mag}., Vol. 30, pp. 23-34 (2009).

\bibitem{ref:Bommasani21}
Bommasani, R., Hudson, D. A., Adeli, E., Altman, R., Arora, S., Arx, von S., Bernstein, M. S., Bohg, J., Bosselut, A., Brunskill, E., et al.: On the opportunities and risks of foundation models, \textit{arXiv preprint arXiv:2108.07258} (2021).

\bibitem{ref:Frich19}
Frich, J., Remy, C., Biskjaer, M. M., Vermeulen, M. L., and Dalsgaard, P.: Mapping the Landscape of Creativity Support Tools in HCI, In Proceedings of the 2019 CHI Conference on Human Factors in Computing Systems, pp. 1-18, (2019). https://doi.org/10.1145/3290605.3300619

\bibitem{ref:OpenAI23}
Ecoffet, A.: GPT-4 Technical Report, \textit{arXiv preprint arXiv:2303.08774v3} (2023).

\bibitem{ref:Zhao23}
Zhao, W. X., Zhou, K., Li, J., Tang, T., Wang, X., Hou, Y., Min, Y., Zhang, B., Zhang, J., Dong, Z., et al.: A Survey of Large Language Models, \textit{arXiv preprint, arXiv.2303.18223} (2023).

\bibitem{ref:Mirowski23}
Mirowski, P., Mathewson, K. W., Pittman, J., and Evans, R.: Co-Writing Screenplays and Theatre Scripts with Language Models: Evaluation by Industry Professionals. In Proceedings of the 2023 CHI Conference on Human Factors in Computing Systems, No. 355, pp. 1–34, (2023).

\bibitem{ref:Chen23}
Chen, B., Zhang, Z., Langrené, N., and Zhu, S.: Unleashing the potential of prompt engineering in Large Language Models: A comprehensive review, \textit{arXiv preprint arXiv.2310.14735} (2023).

\bibitem{ref:Ko23}
Ko, H.-K., Park, G., Jeon, H., Jo, J., Kim, J., Seo, J.: Large-scale Text-to-Image Generation Models for Visual Artists’ Creative Works, In Proceedings of the 28th International Conference on Intelligent User Interfaces, pp.919-93, (2023).

\bibitem{ref:Louie20}
Louie, R., Coenen, A., Huang, C. Z., Terry, M., Cai, C. J.: Novice-AI Music Co-Creation via AI-Steering Tools for Deep Generative Models, In Proceedings of the 2020 CHI Conference on Human Factors in Computing Systems, pp.1-13 (2020).

%\bibitem{ref:Vaswani17}
%Vaswani, A., Shazeer, N., Parmar, N., Uszkoreit, J., Jones, L., Gomez, A. N., Kaiser, Ł., Polosukhin, I.: Attention is All you Need, In Advances in Neural Information Processing Systems 30, (2017).

%\bibitem{ref:Devlin18}
%Devlin, J., Chang, M.-W., Lee, K., Toutanova, K.: BERT: Pre-training of Deep Bidirectional Transformers for Language Understanding, \textit{arXiv preprint arXiv:1810.04805} (2018).
\bibitem{ref:Guzdial19}
Guzdial, M., Liao, N., Chen, J., Chen, S.-Y., Shah, S., Shah, V., Reno, J., Smith, G., Riedl, M. O.: Friend, Collaborator, Student, Manager: How Design of an AI-Driven Game Level Editor Affects Creators, In Proceedings of the 2019 CHI Conference on Human Factors in Computing Systems, No. 624, pp.1-13, (2019).

\bibitem{ref:Osone21}
Osone, H., Lu, J.-L., and Ochiai, Y.: BunCho: AI Supported Story Co-Creation via Unsupervised Multitask Learning to Increase Writers’ Creativity in Japanese, In Abstracts of the 2021 CHI Conference on Human Factors in Computing Systems, No. 19, pp. 1-10,  (2021). 


\bibitem{ref:Radford20}
Radford, A., Wu, J., Child, R., Luan, D., Amodei, D., and Sutskever, I.: Language models are unsupervised multitask learners, \textit{OpenAI blog}, (2019).

\bibitem{ref:Yuan22}
Yuan, A., Coenen, A., Reif, E., and Ippolito, D.: Wordcraft: Story Writing With Large Language Models, In Proceedings of the 27th International Conference on Intelligent User Interfaces, pp. 841-852, (2022).

\bibitem{ref:Thoppilan22}
Thoppilan, R., De Freitas, D., Hall, J., Shazeer, N., Kulshreshtha, A., Cheng, H.-T., Jin, A., Bos, T., Baker, L., Du, Y., et al.: LaMDA: Language Models for Dialog Applications, \textit{arXiv preprint arXiv:2201.08239.}, (2022).

\bibitem{ref:Higasa23}
Hisaga, T., Kawamura, T., Kurihara, S.: Automated Story Generation Based on Narrative Structure to Inspire Creativity, 人工知能学会全国大会論文集, (2023).

\bibitem{ref:Brown20}
Brown, T. B., Mann, B., Ryder, N., Subbiah, M., Kaplan, J., Dhariwal, P., Neelakantan, A., Shyam, P., Sastry, G., Askell, A., Agarwal, S., et al.: Language Models are Few-Shot Learners, \textit{arXiv preprint arXiv:2005.14165}, (2020). 

\bibitem{ref:Tezuka77}
手塚 治虫:ブラック・ジャック， 手塚治虫漫画全集， No.151–168, (1977).

\bibitem{ref:Propp64}
Propp, V. I.: Morphology of the Folktale, Vol. 9, University of Texas Press, (1968).

\bibitem{ref:Jaya11}
Jaya, A. and Uma, G.V.: An Intelligent Automatic Story Generation System by Revising Proppian’s System, CC- SIT 2011, Vol.131, pp.594–603, (2011).

\bibitem{ref:Murai22}
Murai, H.: Automatic Story Structure Generation Based on Fundamental Story Pattern within Existing Stories, 情報処理学会論文誌, Vol. 63, No. 2, pp. 335-346, (2022).

\bibitem{ref:Norva23}
Bakhiyorovna, M. N.: THE SPECIFICATION OF SIGNS AND SYMBOLS IN POETRY, Neo Scientific Peer Reviewed Journal, Vol. 11, pp.43-45, (2023).

\bibitem{ref:Murai21}
Murai, H., Toyosawa, S., Shiratori, T., Yoshida, T., Ishikawa, K., Iwasaki, J., Saito, Y., Nakamura, S., Nemoto, S., Ohta, S., et al.: Extraction of factors that characterize the structures of plot within each story genre, じんもんこん2021論文集, (2021).
% \bibitem{ref:testref1}
% 第１著者, 第２著者, 第３著者: 論文名, 掲載誌名, Vol.~xx, No.~xx, pp.~xxx-xxx, (出版年)

% \bibitem{ref:testref2}
% First A., Second A., and Third A.: Paper Title, Publication Source, Vol.~xx, No.~xx, pp.~xxx-xxx, (year)


\end{thebibliography}

\end{document}
